\documentclass[aspectratio=169]{beamer}

% ===== BEAMER THEMES TEST =====

%\usetheme{AnnArbor}
% \usetheme{Antibes}
% \usetheme{Bergen}
% \usetheme{Berkeley}
% \usetheme{Berlin}
% \usetheme{Boadilla}
% \usetheme{CambridgeUS}
% \usetheme{Copenhagen}
% \usetheme{Darmstadt}
% \usetheme{Dresden}
% \usetheme{Frankfurt}
% \usetheme{Goettingen}
% \usetheme{Hannover}
% \usetheme{Ilmenau}
% \usetheme{JuanLesPins}
% \usetheme{Luebeck}
% \usetheme{Madrid}
% \usetheme{Malmoe}
% \usetheme{Marburg}
% \usetheme{Montpellier}
% \usetheme{PaloAlto}
% \usetheme{Pittsburgh}
% \usetheme{Rochester}
% \usetheme{Singapore}
% \usetheme{Szeged}
 \usetheme{Warsaw}

%\usepackage[faculty=fai]{utbbeamer}
 \usepackage[faculty=utb]{utbbeamer}
% \usepackage[faculty=ft]{utbbeamer}
% \usepackage[faculty=fame]{utbbeamer}
% \usepackage[faculty=fmk]{utbbeamer}
% \usepackage[faculty=fhs]{utbbeamer}
% \usepackage[faculty=flkr]{utbbeamer}

\usepackage{booktabs}
\usepackage{amsmath}

\title[UTB Beamer test]{Testovanie UTB vizuálnej nadstavby pre Beamer}
\subtitle{Ukážkový dokument pre viacero typov snímkov}
\author[Meno Priezvisko]{Meno Priezvisko}
\institute{Univerzita Tomáše Bati ve Zlíně}
\date{\today}

\begin{document}
	
	% -------------------------------------------------
	\begin{frame}
		\titlepage
	\end{frame}
	
	% -------------------------------------------------
	\begin{frame}{Obsah}
		\tableofcontents
	\end{frame}
	
	% =================================================
	\section{Základné prvky}
	
	\begin{frame}{Text a zvýraznenie}
		Toto je bežný text v prezentácii.
		
		\vspace{0.4cm}
		
		Toto je \alert{zvýraznený text pomocou príkazu alert}.
		
		\vspace{0.4cm}
		
		\textbf{Tučný text}, \textit{kurzíva} a \texttt{monospace text}.
	\end{frame}
	
	% -------------------------------------------------
	\begin{frame}{Zoznamy}
		\begin{itemize}
			\item Prvá úroveň zoznamu
			\item Druhá položka
			\begin{itemize}
				\item Vnorená položka
				\item Ďalšia vnorená položka
			\end{itemize}
			\item Posledná položka
		\end{itemize}
		
		\vspace{0.4cm}
		
		\begin{enumerate}
			\item Prvý krok
			\item Druhý krok
			\item Tretí krok
		\end{enumerate}
	\end{frame}
	
	% =================================================
	\section{Bloky}
	
	\begin{frame}{Typy blokov}
		\begin{block}{Bežný blok}
			Toto je bežný informačný blok.
		\end{block}
		
		\begin{alertblock}{Upozornenie}
			Toto je alert blok na dôležité informácie.
		\end{alertblock}
		
		\begin{exampleblock}{Príklad}
			Toto je príkladový blok.
		\end{exampleblock}
	\end{frame}
	
	% -------------------------------------------------
	\begin{frame}{Viac blokov na jednom snímku}
		\begin{columns}
			
			\begin{column}{0.48\textwidth}
				\begin{block}{Výhody}
					\begin{itemize}
						\item jednotná identita,
						\item jednoduché použitie,
						\item kompatibilita s témami.
					\end{itemize}
				\end{block}
			\end{column}
			
			\begin{column}{0.48\textwidth}
				\begin{alertblock}{Riziká}
					\begin{itemize}
						\item rôzne správanie tém,
						\item rozdielne hlavičky,
						\item potreba testovania.
					\end{itemize}
				\end{alertblock}
			\end{column}
			
		\end{columns}
	\end{frame}
	
	% =================================================
	\section{Tabuľky a dáta}
	
	\begin{frame}{Tabuľka}
		\begin{table}
			\centering
			\begin{tabular}{lll}
				\toprule
				Prvok & Beamer príkaz & Účel \\
				\midrule
				Farba & \texttt{setbeamercolor} & vizuálna identita \\
				Font & \texttt{setbeamerfont} & typografia \\
				Pätička & \texttt{setbeamertemplate} & navigácia \\
				Logo & \texttt{includegraphics} & identifikácia \\
				\bottomrule
			\end{tabular}
		\end{table}
	\end{frame}
	
	% -------------------------------------------------
	\begin{frame}{Porovnanie fakúlt}
		\begin{tabular}{ll}
			\toprule
			Voľba balíka & Výsledok \\
			\midrule
			\texttt{faculty=utb} & hlavná UTB identita \\
			\texttt{faculty=fai} & Fakulta aplikované informatiky \\
			\texttt{faculty=ft} & Fakulta technologická \\
			\texttt{faculty=fame} & Fakulta managementu a ekonomiky \\
			\texttt{faculty=fmk} & Fakulta multimediálních komunikací \\
			\texttt{faculty=fhs} & Fakulta humanitních studií \\
			\texttt{faculty=flkr} & Fakulta logistiky a krizového řízení \\
			\bottomrule
		\end{tabular}
	\end{frame}
	
	% =================================================
	\section{Matematika}
	
	\begin{frame}{Matematický zápis}
		Ukážka rovnice:
		
		\[
		E = mc^2
		\]
		
		Ukážka sústavy:
		
		\[
		\begin{aligned}
			a^2 + b^2 &= c^2, \\
			f(x) &= \int_0^x t^2 \, dt.
		\end{aligned}
		\]
		
		\begin{block}{Poznámka}
			Matematika ostáva riadená štandardným LaTeX mechanizmom.
		\end{block}
	\end{frame}
	
	% =================================================
	\section{Obrázky a rozloženie}
	
	\begin{frame}{Logo a text}
		\begin{columns}
			
			\begin{column}{0.45\textwidth}
				\begin{center}
					\includegraphics[width=0.8\textwidth]{\UTBLogoFile}
				\end{center}
			\end{column}
			
			\begin{column}{0.5\textwidth}
				\begin{block}{Použitie loga}
					Logo sa používa na titulnej strane a podľa potreby aj v obsahu prezentácie.
				\end{block}
			\end{column}
			
		\end{columns}
	\end{frame}
	
	% -------------------------------------------------
	\begin{frame}{Dvojstĺpcový layout}
		\begin{columns}[T]
			
			\begin{column}{0.48\textwidth}
				\textbf{Ľavý stĺpec}
				
				\begin{itemize}
					\item analýza identity,
					\item výber farieb,
					\item návrh pravidiel.
				\end{itemize}
			\end{column}
			
			\begin{column}{0.48\textwidth}
				\textbf{Pravý stĺpec}
				
				\begin{enumerate}
					\item implementácia balíka,
					\item testovanie tém,
					\item dokumentácia.
				\end{enumerate}
			\end{column}
			
		\end{columns}
	\end{frame}
	
	% =================================================
	\section{Prekrytia}
	
	\begin{frame}{Postupné zobrazovanie}
		\begin{itemize}
			\item<1-> Najprv sa zobrazí prvý bod.
			\item<2-> Potom druhý bod.
			\item<3-> Nakoniec tretí bod.
		\end{itemize}
		
		\only<4>{
			\begin{alertblock}{Záver}
				Prekrytia fungujú aj po aplikovaní UTB farieb.
			\end{alertblock}
		}
	\end{frame}
	
	% =================================================
	\section{Záver}
	
	\begin{frame}{Zhrnutie testovania}
		\begin{block}{Overené prvky}
			\begin{itemize}
				\item titulná strana,
				\item obsah,
				\item hlavička a pätička,
				\item bloky,
				\item tabuľky,
				\item matematika,
				\item obrázky,
				\item viacstĺpcové rozloženie.
			\end{itemize}
		\end{block}
	\end{frame}
	
	% -------------------------------------------------
	\begin{frame}{Ďakujem za pozornosť}
		\begin{center}
			{\Huge Ďakujem}
			
			\vspace{0.8cm}
			
			\UTBLogoLarge
		\end{center}
	\end{frame}
	
\end{document}